\begin{pastexam}[label={px:ust-115-ee-ds-05}]{臺聯大電機類 115 學年度 第 5 題}
（單選；2 分）

Which of the following situations would lead to Stack Overflow?
\begin{enumerate}[label=(\Alph*)]
\item Try popping data from an empty stack.
\item The recursion is too deep and there is no convergence condition.
\item The front of the circular queue caught up with the rear
\item The number of nodes in the implementation of linked sequences is too small.
\end{enumerate}
\end{pastexam}

\begin{pastsoln}
答案是 (B)。

溢位（overflow）是堆疊已經滿了還要 \proc{Push}（資結書第~6.1~節）。 程式執行時的呼叫堆疊
（call stack）容量有限，每呼叫一次函式就壓一個框（參數、返回位址、區域變數），函式返回才拿掉。
遞迴沒有終止條件（題目說的 convergence condition）或深度太大時，框一直壓上去不會拿掉，最後
超出呼叫堆疊的容量，就是 stack overflow（遞迴和堆疊的關係見資結書第~6.4~節）。

逐項檢查：
\begin{itemize}
\item (A) 對空堆疊 \proc{Pop} 是下溢（underflow），不是溢位。 錯。
\item (B) 對，見上。
\item (C) 這是佇列的空、滿判斷，和堆疊無關。 錯。
\item (D) 「串列的節點太少」不會造成溢位；串列實作的堆疊沒有固定容量，節點少只表示裡面的資料
少。 錯。
\end{itemize}
\end{pastsoln}
